%------------------------------------------------------------------------------%

\usepackage{apresentacao}
\usepackage{mathconfig}
\usepackage{tabto}

\makeatletter
\graphicspath  {{./figs/}}
\def\input@path{{./figs/}}
\makeatother

\newcounter{exemplo}
\setcounter{exemplo}{0}

%------------------------------------------------------------------------------%
\title {Usos da Séries de Taylor -- Arco Tangente}
\author{\large Luis Alberto D'Afonseca}
\date  {Integrais e Séries}

%------------------------------------------------------------------------------%
\begin{document}

%------------------------------------------------------------------------------%
\begin{frame}
    \titlepage
\end{frame}

%------------------------------------------------------------------------------%
\section{Cálculo da Série de Taylor para o Arco Tangente}

%------------------------------------------------------------------------------%
\begin{frame}{Arco Tangente}
  Calcular a Série de Taylor de
  \[
    \arctan(x)
  \]

  \pause
  \vspace{\baselineskip}
  Sabemos que
  \[
    \frac{d}{dx}\arctan(x) = \frac{1}{1+x^2}
  \]
\end{frame}

%------------------------------------------------------------------------------%
\begin{frame}{Arco Tangente}
  Calculando as derivadas
  \begin{align*}
    \action<+->{\frac{d}{dx}     & \arctan(x) = \hpm \frac{1}{1+x^2}               \\[3mm]}
    \action<+->{\frac{d^2}{dx^2} & \arctan(x) =    - \frac{2x}{(1+x^2)^2}          \\[3mm]}
    \action<+->{\frac{d^3}{dx^3} & \arctan(x) = \hpm \frac{6x^2-2}{(1+x^2)^3}      \\[3mm]}
    \action<+->{\frac{d^4}{dx^4} & \arctan(x) =    - \frac{24x(x^2-1)}{(1+x^2)^4}}
  \end{align*}
  \action<+->{\emph{Não é um bom método!}}
\end{frame}

%------------------------------------------------------------------------------%
\begin{frame}{Arco Tangente}
  Cálculos formais (sem rigor)
  \pause
  \[
    \frac{d}{dx} \arctan(x) = \frac{1}{1+x^2}
  \]
  \pause
  Soma de uma série geométrica com $a=1$ e $r=-x^2$
  \[
    \frac{d}{dx} \arctan(x) = \frac{1}{1+x^2}
                            = \frac{a}{1-r}
                            = 1 - x^2 + x^4 - x^6 + x^8 - \cdots
  \]
  \pause
  Integrando os dois lados
  \[
    \arctan(x) = x -\frac{x^3}{3} + \frac{x^5}{5} - \frac{x^7}{7} + \frac{x^9}{9} - \cdots
  \]
\end{frame}

%------------------------------------------------------------------------------%
\begin{frame}{Problema}
  Como sabemos se os limites convergem?
\end{frame}

%------------------------------------------------------------------------------%
\begin{frame}{Demonstrando a Série de Taylor do Arco Tangente}
  \[
    \frac{1}{1+t^2} = 1 - t^2 + t^4 - t^6 + \cdots
                    + (-1)^n t^{2n}
                    + {\color<2->{blue}(-1)^{n+1} t^{2(n+1)} + \cdots}
  \]

  \pause
  \vspace{0.5\baselineskip}
  Usando a fórmula da soma da Série Geométrica a partir de $n+1$
  \[
    \frac{1}{1+t^2} = 1 - t^2 + t^4 - t^6 + \cdots
                    + (-1)^n t^{2n}
                    + {\color{blue}\frac{(-1)^{n+1} t^{2(n+1)}}{1+t^2}}
  \]

  \pause
  \vspace{0.5\baselineskip}
  Agora temos um número finito de termos
\end{frame}

%------------------------------------------------------------------------------%
\begin{frame}{Integrando a Soma Finita}
  Escrevendo a derivada como uma soma finita
  \[
    \frac{d}{dx} \arctan x = 1 - x^2 + x^4 - x^6
                           + \cdots
                           + (-1)^n x^{2n}
                           + \frac{(-1)^{n+1} x^{2(n+1)}}{1+x^2}
  \]

  \pause
  \vspace{0.5\baselineskip}
  Integrando
  \[
    \arctan x = x - \frac{x^3}{3}
                  + \frac{x^5}{5}
                  - \frac{x^7}{7}
                  + \cdots
                  + \frac{(-1)^n x^{2n+1}}{2n+1}
                  + \int_0^x \frac{(-1)^{n+1} t^{2(n+1)}}{1+t^2} \,dt
  \]
\end{frame}

%------------------------------------------------------------------------------%
\begin{frame}{A Série de Taylor Converge se o Erro Tende a Zero}
  \[
    R_n(x) = \int_0^x \frac{(-1)^{n+1} t^{2(n+1)}}{1+t^2} \,dt
  \]
  \pause
  como \(\;1+t^2\geq 1\;\)
  \[
    \abs{R_n(x)} \leq \int_0^{\abs{x}} t^{2(n+1)} \,dt \pause
      \; = \; \frac{t^{2n+3}}{2n+3}\evalat{0}{\abs{x}} \pause
      \; = \; \frac{\abs{x}^{2n+3}}{2n+3}
  \]
  \pause
  Portanto
  \[
    \lim_{n\to\infty} R_n(x) = 0 \qquad\text{quando}\qquad \abs{x} \leq 1
  \]
\end{frame}

%------------------------------------------------------------------------------%
\begin{frame}{Série de Taylor do Arco Tangente}
  Para todo \(\;x\in[-1,\,´1\,]\;\) temos
  \[
    \arctan x = x - \frac{x^3}{3}
                  + \frac{x^5}{5}
                  - \frac{x^7}{7}
                  + \frac{x^9}{9}
                  - \cdots
  \]
\end{frame}

%------------------------------------------------------------------------------%
\begin{frame}{\contentsname}
  \tableofcontents
\end{frame}

%------------------------------------------------------------------------------%
\end{document}
%------------------------------------------------------------------------------%
